\section*{الفصل \ref{ch:g6:water-states} --- الماء في حالاته كلها}

\begin{solution}{exo:g6:water-states:1}
\omterm{def:g6:water-states:changes}{الانصهار} (الصقيع عن عشب الصباح)؛ \omterm{def:g6:water-states:changes}{والتصلّب} (مكعّبات ثلج
المجمّدة)؛ \omterm{def:g6:water-states:changes}{والتبخّر} (جفاف الغسيل، والقِدر الغالي)؛ \omterm{def:g6:water-states:changes}{والتكاثف}
(الغبش على \omterm{def:g5:mirrors-reflection:reflection}{المرآة} الباردة)؛ \omterm{def:g6:water-states:changes}{والتسامي} (كومة ثلج تتقلّص في صقيع
جافّ)؛ ومن الغازية إلى الصلبة مباشرة (ريش الصقيع على الزجاج).
\end{solution}

\begin{solution}{exo:g6:water-states:2}
الندى: \omterm{def:g6:water-states:changes}{تكاثف}. ورجل الثلج المتقلّص بلا بركة: تسامٍ --- من الصلبة
إلى البخار مباشرة في \omterm{def:g2:air-around-us:air}{الهواء} البارد الجافّ. و``دخان'' المجمّدة:
\omterm{def:g6:water-states:changes}{تكاثف} --- إذ \omterm{def:g2:ice-water-steam:evapcond}{يتكاثف} بخار الغرفة غيمة صغيرة في \omterm{def:g2:air-around-us:air}{الهواء} البارد
(قطيرات لا غازًا).
\end{solution}

\begin{solution}{exo:g6:water-states:3}
\qty{412}{g} كما هي: \omterm{def:g3:measuring-mass:mass}{فالكتلة} تنجو من تغيّرات الحالة --- ولم
تفقد القنّينة المحكمة شيئًا ولم تكسب.
\end{solution}

\begin{solution}{exo:g6:water-states:4}
مغلقة، حتى لا ينجرف بخار بغراماته (فالتدقيق يجب أن يبقى
مغلقًا)؛ ومجفّفة، لأن ماءً عالقًا أو صقيعًا من الخارج يضيف
غرامات غريبة تخصّ الغرفة لا التدقيق.
\end{solution}

\begin{solution}{exo:g6:water-states:5}
لم تُفنَ: فقد غادر \qty{200}{g} من الماء بخارًا عبر \omterm{def:g2:air-around-us:air}{هواء}
المطبخ --- \omterm{def:g6:water-states:changes}{والتبخّر} نقل الغرامات خارج القِدر، لا خارج العالم.
\end{solution}

\begin{solution}{exo:g6:water-states:6}
نحو \qty{1.1}{L} من الثلج --- أي حيّزًا أكبر بالعشر. \omterm{def:g3:measuring-mass:mass}{والكتلة} لم
تتغيّر: \qty{1}{kg} داخلة و\qty{1}{kg} خارجة.
\end{solution}

\begin{solution}{exo:g6:water-states:7}
الماء المحبوس في أنبوب ينتفخ حين \omterm{def:g2:ice-water-steam:meltfreeze}{يتجمّد}، فيستسلم الأنبوب مثل
القنّينة الزجاجية. ويفرّغ البستانيون صنابير الخريف حتى لا يبقى
في الداخل ما ينتفخ.
\end{solution}

\begin{solution}{exo:g6:water-states:8}
الثلج يحشو مادة \omterm{def:g6:mass-and-volume:volume}{لتر} في $1.1$ \omterm{def:g6:mass-and-volume:volume}{لتر} من الحيّز: فهو أخفّ من الماء
للحيّز نفسه --- والأخفّ للحيّز نفسه هو بالضبط ما \omterm{def:g1:floating-and-sinking:words}{يطفو}.
\end{solution}

\begin{solution}{exo:g6:water-states:9}
البحر (\qty{1}{kg}، \omterm{def:g3:solids-liquids-gases:liquid}{سائل}) --- \omterm{def:g6:water-states:changes}{تبخّر} --- بخار (\qty{1}{kg}) ---
\omterm{def:g6:water-states:changes}{تكاثف} --- قطيرات غيم (\qty{1}{kg}) --- \omterm{def:g6:water-states:changes}{تصلّب} --- ثلج متساقط
(\qty{1}{kg}) --- \omterm{def:g6:water-states:changes}{انصهار} --- ماء \omterm{def:g6:water-states:changes}{انصهار} (\qty{1}{kg}) --- نهر
إلى البحر (\qty{1}{kg}). ستّ مراحل \omterm{def:g3:measuring-mass:mass}{وكتلة} واحدة لا تتغيّر.
\end{solution}

\begin{solution}{exo:g6:water-states:10}
طمئنه. فالثلج الطافي يزيح أصلًا \emph{وزنه} من الليموناضة؛
\omterm{def:g6:water-states:changes}{والانصهار} يردّ ذلك الماء بعينه، منكمشًا بالعشر الذي كان قد
انتفخه --- فينسلّ إلى الحيّز الذي كان المكعّب يشغله أصلًا تحت
السطح. فتبقى الكأس ممتلئة إلى الحافة، ولا تفيض قطرة.
\end{solution}

\begin{solution}{exo:g6:water-states:11}
باب الغازية إلى الصلبة: بخار خفيّ من \omterm{def:g2:air-around-us:air}{هواء} الغرفة لمس الزجاج
المتجمّد فبنى ثلجًا مباشرة، متخطّيًا الحالة السائلة. وهو ليس
``تجمّدًا'' لأنه لم يوجد على الزجاج \omterm{def:g3:solids-liquids-gases:liquid}{سائل} قطّ --- بل وصل الماء
من الحالة البخارية مباشرة.
\end{solution}

\begin{solution}{exo:g6:water-states:12}
تفقد البركة حرارتها \omterm{def:g2:air-around-us:air}{لهواء} الشتاء عند سطحها، فيبلغ ماء السطح
\qty{0}{\celsius} أولًا ويتكوّن ثلجه هناك --- \omterm{def:g1:floating-and-sinking:words}{ويطفو}، إذ هو
أخفّ من الماء للحيّز نفسه، فيبقى في الأعلى. والغطاء من الثلج
(ومن الثلج المتساقط فوقه) مسرب بطيء: فهو يعزل الماء تحته،
وحرارته لا تتسرّب صاعدة إلا على مهل. فيبقى الماء العميق فوق
درجة \omterm{def:g2:ice-water-steam:meltfreeze}{التجمّد} طوال الشتاء، وتسبح الأسماك تحت سقف بناه لها ماؤها
العجيب المنتفخ.
\end{solution}

\begin{solution}{pb:g6:water-states:1}
\textbf{1.} \omterm{def:g6:water-states:changes}{التبخّر} (\omterm{def:g6:water-states:changes}{تبخّر} سطحي من السطح الدافئ).
\textbf{2.} \qty{1}{kg} كما هي --- \omterm{def:g3:measuring-mass:mass}{فالكتلة} تنجو من تغيّرات
الحالة.
\textbf{3.} \omterm{def:g6:water-states:changes}{التكاثف}.
\textbf{4.} عند \qty{0}{\celsius}.
\textbf{5.} نحو \qty{1.1}{L}.
\textbf{6.} \omterm{def:g3:measuring-mass:mass}{الكتلة} \qty{1}{kg}؛ \omterm{def:g6:mass-and-volume:volume}{والحجم} عائد إلى \qty{1}{L}.
\textbf{7.} $1000 - 250 = \qty{750}{g}$ تواصل في الجدول.
\textbf{8.} لا: فالمقدار \qty{250}{g} يركب \omterm{def:g2:air-around-us:air}{الهواء} بخارًا سليمًا؛
وأرجح أبوابها التالية \omterm{def:g6:water-states:changes}{التكاثف} --- إلى غيم، ثم مطر.
\textbf{9.} $750 + 250 = \qty{1000}{g}$: فقد عاد \omterm{def:g3:measuring-mass:units}{الكيلوغرام}
كله.
\textbf{10.} خفيّ مرتين: بخارًا بعد مغادرة البحر مباشرة،
وبخارًا من جديد فوق بركة الشلّال --- وكلتاهما في الحالة
الغازية.
\textbf{11.} الانتفاش أمر \emph{\omterm{def:g6:mass-and-volume:volume}{الحجم}}: فالثلج المتساقط بلّورات
جليد مشبوكة \omterm{def:g2:air-around-us:air}{بالهواء}، تشغل حيّزًا سخيًّا. أما \emph{\omterm{def:g3:measuring-mass:mass}{كتلة}} مائنا
المعلَّم فبقيت \qty{1}{kg} بالضبط --- والخفّة بالنسبة إلى
الحيّز ليست خفّة في المادة: والزوج الصحيح من الكلمات هو \omterm{def:g3:measuring-mass:mass}{الكتلة}
\omterm{def:g6:mass-and-volume:volume}{والحجم}.
\textbf{12.} \omterm{def:g6:water-states:changes}{تبخّر}، \omterm{def:g6:water-states:changes}{فتكاثف}، \omterm{def:g6:water-states:changes}{فتصلّب}، \omterm{def:g6:water-states:changes}{فانصهار} --- وللرُّبع
المتبخّر: \omterm{def:g6:water-states:changes}{تبخّر}، \omterm{def:g6:water-states:changes}{فتكاثف} (مطرًا)، فالتحاق \omterm{def:g3:solids-liquids-gases:liquid}{بالسائل}. أربعة أبواب
مستعملة؛ ولم يُحتَج في هذه الجولة إلى \omterm{def:g6:water-states:changes}{التسامي} ولا إلى باب
الغازية إلى الصلبة.
\end{solution}
